% =============================================================
%  Metodologia (2,50 pontos)
% =============================================================
\section*{Metodologia}

\textbf{Abordagem.} Pesquisa aplicada, de natureza exploratória-experimental, orientada à análise comparativa de algoritmos (detecção/rastreamento de objetos, reconhecimento de eventos, modelos multimodais de linguagem) combinados a um modelo de linguagem, avaliados como apoio às etapas reais do trabalho do perito de vídeo - não apenas ao cálculo de uma grandeza isolada.

\textbf{Etapas propostas:}

\begin{enumerate}[leftmargin=1.5em]
    \item \textit{Revisão de literatura sobre o trabalho do perito de vídeo}: mapear, a partir de publicações técnicas e acadêmicas da área forense, como é composto esse trabalho (etapas, ferramentas atualmente usadas, critérios técnicos, principais dificuldades).
    \item \textit{Revisão dos algoritmos aplicáveis}, organizando-os por etapa do fluxo pericial que podem apoiar (triagem/localização de trechos, detecção e rastreamento, reconhecimento de eventos, geração de descrição/análise).
    \item \textit{Definição dos critérios de comparação}: acurácia/qualidade da saída, tempo de processamento, custo computacional, e utilidade percebida da análise gerada para compor um laudo.
    \item \textit{Montagem de um conjunto de vídeos de teste} representativos de câmeras de monitoramento urbano (cenários, resoluções e condições variadas), com anotação de referência quando necessário.
    \item \textit{Implementação e testes} de diferentes combinações de algoritmos sobre esse conjunto, registrando desempenho em cada critério definido.
    \item \textit{Análise comparativa dos resultados} e discussão da viabilidade de aplicação na rotina pericial, incluindo limitações técnicas e considerações sobre validade probatória.
\end{enumerate}

\textbf{Métricas de avaliação:} acurácia/precisão das detecções e classificações, tempo de processamento por vídeo, custo computacional, e avaliação qualitativa da utilidade da análise gerada (por exemplo, por peritos que atuem como avaliadores).

\textbf{Produtos e resultados esperados:} caracterização, a partir da literatura, do fluxo de trabalho real do perito de vídeo em câmeras de monitoramento urbano; análise comparativa dos algoritmos testados quanto à acurácia, custo computacional e utilidade para o trabalho pericial; protótipo/arranjo de referência da combinação mais adequada, como produto técnico do mestrado profissional; discussão sobre limitações técnicas, éticas e de validade probatória; e publicação científica descrevendo os resultados.
